%% =========================================================================================
%% TEMPORARY SECTION -- FOR THOMAS'S REVIEW ONLY, TO BE DELETED BEFORE CIRCULATION.
%% Included from main_final.tex via %% =========================================================================================
%% TEMPORARY SECTION -- FOR THOMAS'S REVIEW ONLY, TO BE DELETED BEFORE CIRCULATION.
%% Included from main_final.tex via %% =========================================================================================
%% TEMPORARY SECTION -- FOR THOMAS'S REVIEW ONLY, TO BE DELETED BEFORE CIRCULATION.
%% Included from main_final.tex via %% =========================================================================================
%% TEMPORARY SECTION -- FOR THOMAS'S REVIEW ONLY, TO BE DELETED BEFORE CIRCULATION.
%% Included from main_final.tex via \input{TEMP_timing_changes}; delete that line and this file.
%% Documents what changed when the Altavilla TIMING factor was added to the surprise measure
%% (2026-07-22). Numbers from Results/diagnostics/compare_timing.R, comparing
%%   Results/main_3factor.rda  (Target + FG + QE)  vs  Results/main.rda  (all four factors).
%% =========================================================================================
\clearpage
\section*{\color{blue}[TEMP] Changes when the Timing factor is included}
\addcontentsline{toc}{section}{[TEMP] Changes when the Timing factor is included}
{\color{blue}

\subsection*{What was wrong}

\autoref{sec:measuring} states that the extracted factors ``consist of the target, the timing, the
forward guidance, and the quantitative easing factor'' and that ``by taking the sum of these
factors'' we obtain the surprise measure. The code did not do this: it summed only Target, FG and
QE. The Timing factor \emph{is} built by the construction script
(\texttt{daily\_data\_update2026.R}, as \texttt{EA\_Timing\_shk}) but was dropped from the weekly
\texttt{MPEA} column selection, so it never reached the estimation. Note that the naming is easy to
confuse: the \emph{Target} factor was always included; the missing one is \emph{Timing}.

The factor survives in the weekly \texttt{rawEA} frame, so it could be added exactly, with no
rebuild of the dataset and therefore no risk that intervening data revisions contaminate the
comparison. The composite now sums all four factors, in both
\texttt{STVARestimation.R} and \texttt{univariate.R}. Specification \texttt{ab3} in the registry
reproduces the old three-factor sum.

Timing is a material component: it has a standard deviation of $2.12$ basis points against $5.25$
for the three-factor sum, and the correlation between the three- and four-factor composites is
$0.93$ (for comparison, the change of factor vintage that the paper treats as immaterial had a
correlation above $0.98$). Adding it raises the standard deviation of the surprise from $5.25$ to
$5.69$ basis points.

\subsection*{What did not change}

\begin{itemize}
\item \textbf{The two outlier events are identical} (17 March 2023 and 4 July 2008 remain the
      maximum and minimum surprise), so the exclusion decision and its justification stand.
\item \textbf{The state allocation is unchanged} ($S_t$ correlation $1.000000$): the threshold
      variable does not involve the surprise. \autoref{fig:stallo} is unaffected.
\item \textbf{The pass-through results of \autoref{sec:passthrough} stand}: the one- to three-year
      nominal coefficient remains about half a basis point per basis point at $t\approx 3$
      ($0.51$, $0.49$, $0.41$ at one, two and three years), still declining with maturity and
      insignificant beyond ten years; swap responses remain small and imprecise. No wording in
      \autoref{sec:passthrough} needed changing.
\item \textbf{The overall impulse responses barely move}: correlation of posterior medians across
      all maturities, stances and horizons is $0.985$, with a maximum deviation of $0.45$ basis
      points.
\item \textbf{Every qualitative conclusion of \autoref{sec:results} is preserved} --- the rotation
      of the pass-through curve, the absence of short-end state dependence, the real-rate /
      inflation-compensation decomposition, and the persistence result.
\end{itemize}

\subsection*{What changed, and where the text was updated}

\begin{itemize}
\item \textbf{The long-end state dependence strengthens slightly.} The loose-minus-tight gap at
      thirty years rises from $1.78$ to $2.12$ basis points, with the posterior probability of a
      positive gap rising from $0.983$ to $0.992$; at fifteen years from $1.54$ to $1.72$
      ($0.981 \to 0.986$). At ten years it is unchanged ($1.31 \to 1.32$).
\item \textbf{The nominal response loses significance slightly earlier along the stance
      distribution}: at the sixty-fifth rather than seventieth percentile at ten years, sixtieth
      rather than sixty-fifth at fifteen, and fifty-fifth rather than sixtieth at thirty. Updated
      in \autoref{sec:results}.
\item \textbf{The five-year regime difference is no longer significant.} Previously the difference
      was significant from five years upwards; it now becomes significant only from seven years.
      This was corrected in the discussion of \autoref{fig:IRFylds}.
\item \textbf{The accommodative-Fed response is now almost perfectly flat} across the curve ($1.3$
      basis points at one year and close to $1.5$ at every maturity from three to thirty), where
      before it had a mild hump at three to five years. If anything this sharpens the
      ``rotation, not shift'' framing.
\item \textbf{The short-end reversal is more pronounced.} Under a restrictive Fed the one-year
      response is now $2.0$ basis points against $1.3$ under an accommodative one (previously
      $1.6$ against $1.4$); the one-year gap moves from $-0.25$ to $-0.69$ basis points. It remains
      statistically insignificant, so the text still reports no short-end state dependence rather
      than a reversal.
\item \textbf{Inflation-compensation responses attenuate somewhat} at the tight-Fed extreme
      ($-1.9$ instead of $-2.3$ basis points at one year), while the one-year real-yield response
      is essentially unchanged ($3.91$ against $3.88$). Numbers updated in \autoref{sec:results}.
\end{itemize}

\subsection*{Assessment}

The correction removes a genuine code-versus-text discrepancy that a replicator would have found,
and it does so without disturbing any conclusion. The main results are, if anything, marginally
cleaner: a flatter accommodative-Fed profile and a slightly stronger long-end gap. The one
substantive edit is the five-year significance boundary in the discussion of
\autoref{fig:IRFylds}.

}
\clearpage
; delete that line and this file.
%% Documents what changed when the Altavilla TIMING factor was added to the surprise measure
%% (2026-07-22). Numbers from Results/diagnostics/compare_timing.R, comparing
%%   Results/main_3factor.rda  (Target + FG + QE)  vs  Results/main.rda  (all four factors).
%% =========================================================================================
\clearpage
\section*{\color{blue}[TEMP] Changes when the Timing factor is included}
\addcontentsline{toc}{section}{[TEMP] Changes when the Timing factor is included}
{\color{blue}

\subsection*{What was wrong}

\autoref{sec:measuring} states that the extracted factors ``consist of the target, the timing, the
forward guidance, and the quantitative easing factor'' and that ``by taking the sum of these
factors'' we obtain the surprise measure. The code did not do this: it summed only Target, FG and
QE. The Timing factor \emph{is} built by the construction script
(\texttt{daily\_data\_update2026.R}, as \texttt{EA\_Timing\_shk}) but was dropped from the weekly
\texttt{MPEA} column selection, so it never reached the estimation. Note that the naming is easy to
confuse: the \emph{Target} factor was always included; the missing one is \emph{Timing}.

The factor survives in the weekly \texttt{rawEA} frame, so it could be added exactly, with no
rebuild of the dataset and therefore no risk that intervening data revisions contaminate the
comparison. The composite now sums all four factors, in both
\texttt{STVARestimation.R} and \texttt{univariate.R}. Specification \texttt{ab3} in the registry
reproduces the old three-factor sum.

Timing is a material component: it has a standard deviation of $2.12$ basis points against $5.25$
for the three-factor sum, and the correlation between the three- and four-factor composites is
$0.93$ (for comparison, the change of factor vintage that the paper treats as immaterial had a
correlation above $0.98$). Adding it raises the standard deviation of the surprise from $5.25$ to
$5.69$ basis points.

\subsection*{What did not change}

\begin{itemize}
\item \textbf{The two outlier events are identical} (17 March 2023 and 4 July 2008 remain the
      maximum and minimum surprise), so the exclusion decision and its justification stand.
\item \textbf{The state allocation is unchanged} ($S_t$ correlation $1.000000$): the threshold
      variable does not involve the surprise. \autoref{fig:stallo} is unaffected.
\item \textbf{The pass-through results of \autoref{sec:passthrough} stand}: the one- to three-year
      nominal coefficient remains about half a basis point per basis point at $t\approx 3$
      ($0.51$, $0.49$, $0.41$ at one, two and three years), still declining with maturity and
      insignificant beyond ten years; swap responses remain small and imprecise. No wording in
      \autoref{sec:passthrough} needed changing.
\item \textbf{The overall impulse responses barely move}: correlation of posterior medians across
      all maturities, stances and horizons is $0.985$, with a maximum deviation of $0.45$ basis
      points.
\item \textbf{Every qualitative conclusion of \autoref{sec:results} is preserved} --- the rotation
      of the pass-through curve, the absence of short-end state dependence, the real-rate /
      inflation-compensation decomposition, and the persistence result.
\end{itemize}

\subsection*{What changed, and where the text was updated}

\begin{itemize}
\item \textbf{The long-end state dependence strengthens slightly.} The loose-minus-tight gap at
      thirty years rises from $1.78$ to $2.12$ basis points, with the posterior probability of a
      positive gap rising from $0.983$ to $0.992$; at fifteen years from $1.54$ to $1.72$
      ($0.981 \to 0.986$). At ten years it is unchanged ($1.31 \to 1.32$).
\item \textbf{The nominal response loses significance slightly earlier along the stance
      distribution}: at the sixty-fifth rather than seventieth percentile at ten years, sixtieth
      rather than sixty-fifth at fifteen, and fifty-fifth rather than sixtieth at thirty. Updated
      in \autoref{sec:results}.
\item \textbf{The five-year regime difference is no longer significant.} Previously the difference
      was significant from five years upwards; it now becomes significant only from seven years.
      This was corrected in the discussion of \autoref{fig:IRFylds}.
\item \textbf{The accommodative-Fed response is now almost perfectly flat} across the curve ($1.3$
      basis points at one year and close to $1.5$ at every maturity from three to thirty), where
      before it had a mild hump at three to five years. If anything this sharpens the
      ``rotation, not shift'' framing.
\item \textbf{The short-end reversal is more pronounced.} Under a restrictive Fed the one-year
      response is now $2.0$ basis points against $1.3$ under an accommodative one (previously
      $1.6$ against $1.4$); the one-year gap moves from $-0.25$ to $-0.69$ basis points. It remains
      statistically insignificant, so the text still reports no short-end state dependence rather
      than a reversal.
\item \textbf{Inflation-compensation responses attenuate somewhat} at the tight-Fed extreme
      ($-1.9$ instead of $-2.3$ basis points at one year), while the one-year real-yield response
      is essentially unchanged ($3.91$ against $3.88$). Numbers updated in \autoref{sec:results}.
\end{itemize}

\subsection*{Assessment}

The correction removes a genuine code-versus-text discrepancy that a replicator would have found,
and it does so without disturbing any conclusion. The main results are, if anything, marginally
cleaner: a flatter accommodative-Fed profile and a slightly stronger long-end gap. The one
substantive edit is the five-year significance boundary in the discussion of
\autoref{fig:IRFylds}.

}
\clearpage
; delete that line and this file.
%% Documents what changed when the Altavilla TIMING factor was added to the surprise measure
%% (2026-07-22). Numbers from Results/diagnostics/compare_timing.R, comparing
%%   Results/main_3factor.rda  (Target + FG + QE)  vs  Results/main.rda  (all four factors).
%% =========================================================================================
\clearpage
\section*{\color{blue}[TEMP] Changes when the Timing factor is included}
\addcontentsline{toc}{section}{[TEMP] Changes when the Timing factor is included}
{\color{blue}

\subsection*{What was wrong}

\autoref{sec:measuring} states that the extracted factors ``consist of the target, the timing, the
forward guidance, and the quantitative easing factor'' and that ``by taking the sum of these
factors'' we obtain the surprise measure. The code did not do this: it summed only Target, FG and
QE. The Timing factor \emph{is} built by the construction script
(\texttt{daily\_data\_update2026.R}, as \texttt{EA\_Timing\_shk}) but was dropped from the weekly
\texttt{MPEA} column selection, so it never reached the estimation. Note that the naming is easy to
confuse: the \emph{Target} factor was always included; the missing one is \emph{Timing}.

The factor survives in the weekly \texttt{rawEA} frame, so it could be added exactly, with no
rebuild of the dataset and therefore no risk that intervening data revisions contaminate the
comparison. The composite now sums all four factors, in both
\texttt{STVARestimation.R} and \texttt{univariate.R}. Specification \texttt{ab3} in the registry
reproduces the old three-factor sum.

Timing is a material component: it has a standard deviation of $2.12$ basis points against $5.25$
for the three-factor sum, and the correlation between the three- and four-factor composites is
$0.93$ (for comparison, the change of factor vintage that the paper treats as immaterial had a
correlation above $0.98$). Adding it raises the standard deviation of the surprise from $5.25$ to
$5.69$ basis points.

\subsection*{What did not change}

\begin{itemize}
\item \textbf{The two outlier events are identical} (17 March 2023 and 4 July 2008 remain the
      maximum and minimum surprise), so the exclusion decision and its justification stand.
\item \textbf{The state allocation is unchanged} ($S_t$ correlation $1.000000$): the threshold
      variable does not involve the surprise. \autoref{fig:stallo} is unaffected.
\item \textbf{The pass-through results of \autoref{sec:passthrough} stand}: the one- to three-year
      nominal coefficient remains about half a basis point per basis point at $t\approx 3$
      ($0.51$, $0.49$, $0.41$ at one, two and three years), still declining with maturity and
      insignificant beyond ten years; swap responses remain small and imprecise. No wording in
      \autoref{sec:passthrough} needed changing.
\item \textbf{The overall impulse responses barely move}: correlation of posterior medians across
      all maturities, stances and horizons is $0.985$, with a maximum deviation of $0.45$ basis
      points.
\item \textbf{Every qualitative conclusion of \autoref{sec:results} is preserved} --- the rotation
      of the pass-through curve, the absence of short-end state dependence, the real-rate /
      inflation-compensation decomposition, and the persistence result.
\end{itemize}

\subsection*{What changed, and where the text was updated}

\begin{itemize}
\item \textbf{The long-end state dependence strengthens slightly.} The loose-minus-tight gap at
      thirty years rises from $1.78$ to $2.12$ basis points, with the posterior probability of a
      positive gap rising from $0.983$ to $0.992$; at fifteen years from $1.54$ to $1.72$
      ($0.981 \to 0.986$). At ten years it is unchanged ($1.31 \to 1.32$).
\item \textbf{The nominal response loses significance slightly earlier along the stance
      distribution}: at the sixty-fifth rather than seventieth percentile at ten years, sixtieth
      rather than sixty-fifth at fifteen, and fifty-fifth rather than sixtieth at thirty. Updated
      in \autoref{sec:results}.
\item \textbf{The five-year regime difference is no longer significant.} Previously the difference
      was significant from five years upwards; it now becomes significant only from seven years.
      This was corrected in the discussion of \autoref{fig:IRFylds}.
\item \textbf{The accommodative-Fed response is now almost perfectly flat} across the curve ($1.3$
      basis points at one year and close to $1.5$ at every maturity from three to thirty), where
      before it had a mild hump at three to five years. If anything this sharpens the
      ``rotation, not shift'' framing.
\item \textbf{The short-end reversal is more pronounced.} Under a restrictive Fed the one-year
      response is now $2.0$ basis points against $1.3$ under an accommodative one (previously
      $1.6$ against $1.4$); the one-year gap moves from $-0.25$ to $-0.69$ basis points. It remains
      statistically insignificant, so the text still reports no short-end state dependence rather
      than a reversal.
\item \textbf{Inflation-compensation responses attenuate somewhat} at the tight-Fed extreme
      ($-1.9$ instead of $-2.3$ basis points at one year), while the one-year real-yield response
      is essentially unchanged ($3.91$ against $3.88$). Numbers updated in \autoref{sec:results}.
\end{itemize}

\subsection*{Assessment}

The correction removes a genuine code-versus-text discrepancy that a replicator would have found,
and it does so without disturbing any conclusion. The main results are, if anything, marginally
cleaner: a flatter accommodative-Fed profile and a slightly stronger long-end gap. The one
substantive edit is the five-year significance boundary in the discussion of
\autoref{fig:IRFylds}.

}
\clearpage
; delete that line and this file.
%% Documents what changed when the Altavilla TIMING factor was added to the surprise measure
%% (2026-07-22). Numbers from Results/diagnostics/compare_timing.R, comparing
%%   Results/main_3factor.rda  (Target + FG + QE)  vs  Results/main.rda  (all four factors).
%% =========================================================================================
\clearpage
\section*{\color{blue}[TEMP] Changes when the Timing factor is included}
\addcontentsline{toc}{section}{[TEMP] Changes when the Timing factor is included}
{\color{blue}

\subsection*{What was wrong}

\autoref{sec:measuring} states that the extracted factors ``consist of the target, the timing, the
forward guidance, and the quantitative easing factor'' and that ``by taking the sum of these
factors'' we obtain the surprise measure. The code did not do this: it summed only Target, FG and
QE. The Timing factor \emph{is} built by the construction script
(\texttt{daily\_data\_update2026.R}, as \texttt{EA\_Timing\_shk}) but was dropped from the weekly
\texttt{MPEA} column selection, so it never reached the estimation. Note that the naming is easy to
confuse: the \emph{Target} factor was always included; the missing one is \emph{Timing}.

The factor survives in the weekly \texttt{rawEA} frame, so it could be added exactly, with no
rebuild of the dataset and therefore no risk that intervening data revisions contaminate the
comparison. The composite now sums all four factors, in both
\texttt{STVARestimation.R} and \texttt{univariate.R}. Specification \texttt{ab3} in the registry
reproduces the old three-factor sum.

Timing is a material component: it has a standard deviation of $2.12$ basis points against $5.25$
for the three-factor sum, and the correlation between the three- and four-factor composites is
$0.93$ (for comparison, the change of factor vintage that the paper treats as immaterial had a
correlation above $0.98$). Adding it raises the standard deviation of the surprise from $5.25$ to
$5.69$ basis points.

\subsection*{What did not change}

\begin{itemize}
\item \textbf{The two outlier events are identical} (17 March 2023 and 4 July 2008 remain the
      maximum and minimum surprise), so the exclusion decision and its justification stand.
\item \textbf{The state allocation is unchanged} ($S_t$ correlation $1.000000$): the threshold
      variable does not involve the surprise. \autoref{fig:stallo} is unaffected.
\item \textbf{The pass-through results of \autoref{sec:passthrough} stand}: the one- to three-year
      nominal coefficient remains about half a basis point per basis point at $t\approx 3$
      ($0.51$, $0.49$, $0.41$ at one, two and three years), still declining with maturity and
      insignificant beyond ten years; swap responses remain small and imprecise. No wording in
      \autoref{sec:passthrough} needed changing.
\item \textbf{The overall impulse responses barely move}: correlation of posterior medians across
      all maturities, stances and horizons is $0.985$, with a maximum deviation of $0.45$ basis
      points.
\item \textbf{Every qualitative conclusion of \autoref{sec:results} is preserved} --- the rotation
      of the pass-through curve, the absence of short-end state dependence, the real-rate /
      inflation-compensation decomposition, and the persistence result.
\end{itemize}

\subsection*{What changed, and where the text was updated}

\begin{itemize}
\item \textbf{The long-end state dependence strengthens slightly.} The loose-minus-tight gap at
      thirty years rises from $1.78$ to $2.12$ basis points, with the posterior probability of a
      positive gap rising from $0.983$ to $0.992$; at fifteen years from $1.54$ to $1.72$
      ($0.981 \to 0.986$). At ten years it is unchanged ($1.31 \to 1.32$).
\item \textbf{The nominal response loses significance slightly earlier along the stance
      distribution}: at the sixty-fifth rather than seventieth percentile at ten years, sixtieth
      rather than sixty-fifth at fifteen, and fifty-fifth rather than sixtieth at thirty. Updated
      in \autoref{sec:results}.
\item \textbf{The five-year regime difference is no longer significant.} Previously the difference
      was significant from five years upwards; it now becomes significant only from seven years.
      This was corrected in the discussion of \autoref{fig:IRFylds}.
\item \textbf{The accommodative-Fed response is now almost perfectly flat} across the curve ($1.3$
      basis points at one year and close to $1.5$ at every maturity from three to thirty), where
      before it had a mild hump at three to five years. If anything this sharpens the
      ``rotation, not shift'' framing.
\item \textbf{The short-end reversal is more pronounced.} Under a restrictive Fed the one-year
      response is now $2.0$ basis points against $1.3$ under an accommodative one (previously
      $1.6$ against $1.4$); the one-year gap moves from $-0.25$ to $-0.69$ basis points. It remains
      statistically insignificant, so the text still reports no short-end state dependence rather
      than a reversal.
\item \textbf{Inflation-compensation responses attenuate somewhat} at the tight-Fed extreme
      ($-1.9$ instead of $-2.3$ basis points at one year), while the one-year real-yield response
      is essentially unchanged ($3.91$ against $3.88$). Numbers updated in \autoref{sec:results}.
\end{itemize}

\subsection*{Assessment}

The correction removes a genuine code-versus-text discrepancy that a replicator would have found,
and it does so without disturbing any conclusion. The main results are, if anything, marginally
cleaner: a flatter accommodative-Fed profile and a slightly stronger long-end gap. The one
substantive edit is the five-year significance boundary in the discussion of
\autoref{fig:IRFylds}.

}
\clearpage
